\section{Investigación publicada en 2026}
La revisión incluye nueve papers enviados inicialmente a arXiv entre febrero y septiembre de 2026. Se comprobaron las fechas de envío por separado de las revisiones; las versiones consultadas figuran en las referencias. La selección se centró en coordinación, traspaso de contexto, asignación de skills y eficiencia medible. Es una revisión técnica focalizada, no una revisión sistemática exhaustiva. Sus puntuaciones no se combinan con el experimento posterior.

\subsection{Topología de ejecución y autoridad}
Kulkarni y Kulkarni comparan flujos secuenciales, paralelos, supervisor--trabajador y reflexivos para extraer documentos financieros. Sus descripciones motivan las cuatro familias estudiadas aquí. El dominio y la configuración de modelos/precios de enero de 2025 limitan la extrapolación a un asistente personal; no se reproducen sus configuraciones ni su rendimiento publicado.\source{https://arxiv.org/html/2603.22651v1}{Kulkarni y Kulkarni, arXiv:2603.22651v1, 2026}

DESBench estudia coordinación centralizada, jerárquica, heterárquica y holónica en planificación industrial. Su heterarquía implementada utiliza pujas mediadas. Amplía la pregunta desde la forma del grafo hasta quién puede asignar trabajo, rechazar compromisos y resolver conflictos; su simulador industrial constituye otro entorno de evaluación.\source{https://arxiv.org/html/2605.13172v1}{Wang et al., arXiv:2605.13172v1, 2026}

ParaGUI envía trabajadores de interfaz gráfica en paralelo por rondas e integra sus resúmenes. Aporta evidencia sobre trabajo interactivo divisible e identifica la ejecución del trabajador como una limitación. Sus controles del historial visual y el entorno de varios escritorios impiden trasladar sus cifras directamente a este experimento textual.\source{https://arxiv.org/html/2607.22689v1}{Yu et al., arXiv:2607.22689v1, 2026}

\subsection{Traspaso de información y evaluación}
OrchBench representa dependencias y transferencias entre agentes y simula calidad, tiempo de terminación y tokens. Su mecanismo de retención convierte la pérdida de información en una variable de coordinación. Aquí se ejecutan todos los trabajadores contra la API; no se sustituyen consumo o respuestas observados por resultados o costes simulados.\source{https://arxiv.org/html/2607.25656v1}{Ren et al., arXiv:2607.25656v1, 2026}

CRAFT entrega vistas parciales complementarias de una construcción. Distingue razonamiento individual y comunicación efectiva y describe correcciones repetidas improductivas. Motiva conservar evidencia y evaluar el resultado final. Como todos nuestros trabajadores reciben la tarea original, no se estudia su problema de observación parcial.\source{https://arxiv.org/html/2603.25268v2}{Nath, VanderHoeven y Krishnaswamy, arXiv:2603.25268v2, 2026}

\subsection{Skills y contexto específico}
\textit{Subagents vs Agent Skills} distingue paquetes procedurales con interfaces explícitas de entrada/salida y texto poco estructurado. Su evaluación principal utiliza un subconjunto seleccionado de 64 tareas e informa menor contexto máximo con mayor consumo total. Esto respalda estudiar la carga de información por separado de la topología.\source{https://arxiv.org/html/2609.09233v1}{Piriyakulkij et al., arXiv:2609.09233v1, 2026}

EvoAgent combina delegación jerárquica y selección por palabras clave, embeddings y modelo. Su evaluación de dominio utiliza jueces modelo; esas puntuaciones no establecen fiabilidad de ejecución exacta. Se toma como referencia la separación de recuperación y delegación, sin implementar aprendizaje autónomo ni reproducir su cascada de selección.\source{https://arxiv.org/html/2604.20133v3}{A. Zhang et al., arXiv:2604.20133v3, 2026}

Swarm Skills empaqueta roles, flujos, límites y dependencias como recursos reutilizables de coordinación. Su caso cualitativo deja pendiente validar ampliamente la compatibilidad. Se consideran las descripciones portables una propuesta de diseño, no evidencia de que cualquier framework multiagente pueda intercambiarlas sin comprobar su funcionamiento.\source{https://arxiv.org/html/2605.10052v2}{X. Zhang et al., arXiv:2605.10052v2, 2026}

Xu y Yan revisan la carga progresiva: descubrir metadatos, cargar instrucciones elegidas y recuperar recursos de apoyo. Esto distingue el catálogo del contenido insertado en cada contexto. Como revisión, orienta la política de carga, pero no aporta un resultado de rendimiento independiente para nuestras arquitecturas.\source{https://arxiv.org/html/2602.12430v4}{Xu y Yan, arXiv:2602.12430v4, 2026}
